% =============================================================================
% METHODS
% Restructured (Round 2): Target design → Implemented design → Why gap.
% No "ideal method then withdraw"; target and implemented are separated upfront.
% =============================================================================

\section{Methods}
\label{sec:methods}

\subsection{Target Research Design}
\label{sec:target-design}

To \emph{identify} whether enterprise digitalization improves both TFP and environmental performance (the double dividend) among heavily polluting firms, the \emph{target} design would require the following. \textbf{Constructs:} (1) \emph{Digitalization} from annual-report text (e.g., keyword count or TF--IDF of digital-related terms). (2) \emph{TFP} from a production-function residual (e.g., Levinsohn--Petrin or Olley--Pakes) using output, capital, and labor. (3) \emph{Environmental performance} from emission intensity or green total factor productivity (GTFP). (4) \emph{Green innovation} (mediator) from green patent applications or grants. (5) \emph{Environmental regulation} (moderator) from regional indices (e.g., enforcement, investment). (6) \emph{Size} from log total assets. \textbf{Design:} Firm-year panel; firm and year fixed effects; digitalization lagged one year; standard errors clustered at firm level. For causal identification, instrumental variables (e.g., historical communication infrastructure or peer digitalization) with a strong first stage (e.g., $F > 10$) would be needed. This subsection describes the estimand and constructs we aim to identify; the next subsection describes what we actually implement.

\subsection{Implemented Empirical Design}
\label{sec:implemented-design}

\subsubsection{Sample and data sources}

The sample consists of Chinese A-share listed firms in industries classified as \emph{heavily polluting} (e.g., thermal power, steel, chemicals, cement, nonferrous metals) according to the official Chinese environmental verification list (industry classification codes); selection is not subjective. Data are from open, publicly available sources: financial and accounting data from AKShare (e.g., yjbb income-statement data) and listed firms' public disclosures; balance-sheet data from a separate merge (see below). The panel spans 2019--2022. After standard filters (excluding financials, ST/PT, key missing), we have 3,481 firm-year observations from 878 unique firms with non-missing revenue and ROE. All continuous variables are winsorized at the 1st and 99th percentiles.

\subsubsection{Sample construction, merge diagnostics, and flowchart}

To document the data pipeline and the origin of the $N = 64$ total-assets match, we report merge steps and diagnostics explicitly.

\textbf{Data sources and identifiers.} (1) \emph{Income-statement (profit) data:} Source \texttt{yjbb} via AKShare; firm identifier \texttt{stock\_code} (string); \texttt{year}; variables include \texttt{revenue}, \texttt{roe}, \texttt{gross\_margin}. (2) \emph{Heavy-pollution list:} List of \texttt{stock\_code} for firms in heavily polluting industries; we keep only firm-years whose \texttt{stock\_code} appears in this list. (3) \emph{Balance-sheet data:} Source \texttt{balance\_sheet\_sina} (or equivalent); firm identifier and \texttt{year}; variable \texttt{total\_assets}. Merge key is (\texttt{stock\_code}, \texttt{year}). Alignment is by calendar year; reporting period is assumed annual.

\textbf{Why total assets has only 64 matches.} After merging the income-statement source (which yields 3,481 firm-years with revenue and ROE for our heavy-pollution subset) to the balance-sheet source on (\texttt{stock\_code}, \texttt{year}), only 64 firm-years have non-missing \texttt{total\_assets}. We verified that identifier format and calendar alignment are consistent; the drop is attributable to differing coverage between the two open sources (ticker lists, time windows, or missing-data patterns), not to a correctable merge bug. Commercial integrated databases (e.g., CSMAR, WIND, RESSET) typically provide consistent firm identifiers and broader asset coverage; we use open sources for transparency and reproducibility and report the merge diagnostic for full disclosure. We therefore use size = ln(revenue) in all regressions so that the regression sample is not restricted to the 64 asset-matched observations. The regression sample is 2,598 firm-years (873 firms) after requiring non-missing lagged digitalization proxy and outcome proxies (see below).

Table~\ref{tab:merge-diagnostics} and Figure~\ref{fig:sample-flow} summarize the sample construction. \textbf{Sample construction (steps):} (1) Income statement + heavy-pollution filter $\rightarrow$ 3,481 firm-years (878 firms). (2) Merge to balance-sheet on (\texttt{stock\_code}, \texttt{year}) $\rightarrow$ 64 firm-years with \texttt{total\_assets}. (3) Regression sample: require non-missing lagged digitalization proxy, TFP proxy, environmental performance proxy, and controls $\rightarrow$ 2,598 firm-years (873 firms). All main regression tables use this final sample (2,598 observations, 873 firms).

\begin{table}[htbp]
  \centering
  \caption{Sample construction and merge diagnostics.}
  \label{tab:merge-diagnostics}
  \small
  \begin{tabularx}{\textwidth}{@{}Xrrl@{}}
    \toprule
    Step & Firm-years & Firms & Note \\
    \midrule
    Income-statement (yjbb) + heavy-pollution filter & 3,481 & 878 & Revenue, ROE available \\
    Merge to balance-sheet (stock\_code, year) & 64 & -- & Total assets available \\
    Regression sample (non-missing lag dig, outcomes, controls) & 2,598 & 873 & Size = ln(revenue) \\
    \bottomrule
  \end{tabularx}
  \medskip
  \parbox{\textwidth}{\raggedright\footnotesize\textit{Note.} Merge key: \texttt{stock\_code}, \texttt{year}. The low match count (64) for total assets is attributable to differing coverage or identifier alignment between the income-statement and balance-sheet sources; see text.}
\end{table}

\subsubsection{Implemented variables}

We implement \emph{proxies} that do not measure the target constructs:

\begin{itemize}
  \item \textbf{Digitalization (implemented):} Industry-year revenue rank (percentile within industry $\times$ year), lagged one year. Derived from \texttt{revenue}; conflates size, market position, and industry-year shocks.
  \item \textbf{TFP (implemented):} ROE (return on equity). Accounting profitability, not a production-function residual.
  \item \textbf{Environmental performance (implemented):} Revenue growth (year-on-year change in log revenue). Reflects business expansion, not emissions or green efficiency.
  \item \textbf{Green innovation / mediator (implemented):} ROE improvement (year-on-year change in ROE, non-negative part). Does not measure green patents or green innovation.
  \item \textbf{Environmental regulation / moderator (implemented):} Calendar year (linear trend). Not a measure of regional regulation strength.
  \item \textbf{Size (implemented):} Log revenue. Preferred target is log total assets; total assets are available for only 64 firm-years after the balance-sheet merge (see merge diagnostics), so we use ln(revenue) to retain the full sample.
\end{itemize}

\subsubsection{Implemented equations}

We estimate the same \emph{structure} as the target design (firm and year fixed effects, lagged predictor), but with implemented proxies. Main equation: $Y_{it} = \alpha + \beta X_{i,t-1} + \mathbf{Z}_{it}'\boldsymbol{\gamma} + \mu_i + \lambda_t + \varepsilon_{it}$, where $X_{i,t-1}$ is the digitalization proxy (industry-year revenue rank, lagged one year), $Y_{it}$ is either the TFP proxy (ROE) or the environmental performance proxy (revenue growth), and $\mathbf{Z}_{it}$ is the vector of controls (size = ln revenue, firm and year fixed effects). We also estimate mediation and moderation \emph{equations} using the implemented mediator and moderators; we do not interpret these as tests of the conceptual mediation or moderation (see Section~\ref{sec:why-gap}). Standard errors are clustered at the firm level where applicable.

\subsection{Why the Implemented Design Illustrates Measurement Failure}
\label{sec:why-gap}

The implemented design does not identify the target constructs; it illustrates \emph{measurement failure} for the following reasons.

\textbf{Construct mismatch.} (1) The digitalization proxy (industry-year revenue rank) is not a measure of digital technology adoption; it reflects revenue-based ranking. (2) The TFP proxy (ROE) is an accounting outcome, not a productivity residual. (3) The environmental performance proxy (revenue growth) does not capture emissions, emission intensity, or GTFP. (4) The mediator proxy (ROE improvement) does not measure green innovation. (5) The moderator (year) is a time trend, not environmental regulation strength. (6) Size is ln(revenue) rather than ln(assets) due to data limitations. Therefore, the estimated coefficients in our tables do not identify the effect of digitalization on TFP or environmental performance; they describe associations among proxy variables.

\textbf{Mechanical and same-source concerns.} The regressor (revenue rank), control (ln revenue), and one outcome (revenue growth) all depend on revenue. This creates a risk of mechanical or same-source correlation. The estimated ``effect'' of the digitalization proxy on the environmental performance proxy may reflect revenue-based variation rather than a causal relationship between digitalization and environmental performance. We do not treat the current estimates as evidence of the latter.

\textbf{Implications.} The results in this paper are \emph{not} evidence on the double dividend. They illustrate \emph{measurement failure}: how revenue-based proxies produce associations that cannot support or refute the theory. The value of the paper lies in answering ``why do revenue-based proxies mislead?'' and in providing a blueprint for future work with direct measures (text-based digitalization, LP/OP TFP, emission-based or GTFP environmental performance, green patents, regional regulation). Reproducibility is subject to the documented variable construction and sample limitations; data and code will be deposited in a public repository upon acceptance.

\subsection{Formal characterization and diagnostic for same-source proxy bias}
\label{sec:formal-diagnostic}

To move from critique to a usable methodological contribution, we provide a formal characterization of same-source proxy bias and a \textbf{diagnostic} that applied researchers can use to assess whether their estimates are likely driven by proxy structure rather than by the target causal effect.

\subsubsection{Setup and notation}

Let $D^*$ denote the target construct (e.g., true digitalization) and $Y^*$ the target outcome (e.g., environmental performance or TFP). The researcher is interested in the causal effect of $D^*$ on $Y^*$. Suppose the researcher does not observe $D^*$ and $Y^*$ but instead uses proxies $X$ (e.g., revenue rank) and $Y$ (e.g., revenue growth or ROE). Let $R$ denote revenue (or a vector of revenue-related variables). We say that the regressor and outcome are \emph{same-source} if both $X$ and $Y$ are measurable functions of $R$ (possibly with independent noise). For example, $X = \phi(R)$ (e.g., industry-year rank of $R$) and $Y = \psi(R)$ (e.g., growth rate of $R$ or an accounting return derived from $R$). The researcher runs the regression
\begin{equation}
Y_{it} = \alpha + \beta X_{i,t-1} + \mathbf{Z}_{it}'\boldsymbol{\gamma} + \varepsilon_{it},
\end{equation}
where $\mathbf{Z}$ includes controls (e.g., firm and year fixed effects, size). The OLS estimator $\hat{\beta}$ is the object of interest.

\begin{quote}
\textbf{Proposition 1 (Same-source proxy bias).} \textit{Suppose (i) $X$ and $Y$ are both measurable functions of the same underlying variable $R$ (and possibly orthogonal noise); (ii) the target constructs $D^*$ and $Y^*$ satisfy $\mathbb{E}[Y^* \mid D^*, R] = \mathbb{E}[Y^* \mid R]$ (no causal effect of $D^*$ on $Y^*$ given $R$, or $D^* \perp Y^* \mid R$). Then the population regression coefficient $\beta$ in the regression of $Y$ on $X$ (controlling for $\mathbf{Z}$) does not identify the causal effect of $D^*$ on $Y^*$. It is determined by the covariance structure of $(X,Y)$ induced by their joint dependence on $R$. In particular, $\beta$ can be nonzero and the corresponding $t$-statistic can be large in absolute value even when the true causal effect of $D^*$ on $Y^*$ is zero.}
\end{quote}

\noindent\textit{Proof.} Let $(\Omega, \mathcal{F}, P)$ be the probability space. Assume that $\mathbf{Z}$ is a vector of controls (possibly including constants and fixed effects), and that the population regression of $Y$ on $(1, X, \mathbf{Z})$ is well defined with finite second moments. The population coefficient $\beta$ is the coefficient on $X$ in the linear projection of $Y$ onto $(1, X, \mathbf{Z})$. Let $\tilde{X}$ denote the population residual from the linear projection of $X$ onto $(1, \mathbf{Z})$, i.e., $\tilde{X} = X - L(X \mid 1, \mathbf{Z})$ where $L(\cdot \mid 1, \mathbf{Z})$ denotes the linear projection. Then
\begin{equation}
\beta = \frac{\mathrm{Cov}(Y, \tilde{X})}{\mathrm{Var}(\tilde{X})},
\end{equation}
provided $\mathrm{Var}(\tilde{X}) > 0$.

\textbf{Step 1.} By assumption (i), there exist measurable functions $\phi$ and $\psi$ and random variables $\eta_X$, $\eta_Y$ such that $X = \phi(R) + \eta_X$ and $Y = \psi(R) + \eta_Y$, with $\mathbb{E}[\eta_X \mid R] = \mathbb{E}[\eta_Y \mid R] = 0$ and $\mathbb{E}[\eta_X \eta_Y \mid R] = 0$ (orthogonal noise given $R$). Then $\mathrm{Cov}(Y, X \mid R) = \mathrm{Cov}(\eta_Y, \eta_X \mid R)$; if we assume $\eta_X$ and $\eta_Y$ are uncorrelated given $R$, then $\mathbb{E}[XY \mid R] = \phi(R)\psi(R) + \mathbb{E}[\eta_X\mid R]\psi(R) + \phi(R)\mathbb{E}[\eta_Y\mid R] + \mathbb{E}[\eta_X\eta_Y\mid R] = \phi(R)\psi(R)$. Hence $\mathrm{Cov}(X, Y) = \mathrm{Cov}(\phi(R), \psi(R)) + \mathrm{Cov}(\eta_X, \eta_Y)$. If $\eta_X$ and $\eta_Y$ are uncorrelated (e.g., independent), then $\mathrm{Cov}(X, Y) = \mathrm{Cov}(\phi(R), \psi(R))$, so the covariance between $X$ and $Y$ is entirely determined by the joint distribution of $R$ and the functions $\phi$, $\psi$.

\textbf{Step 2.} The residual $\tilde{X} = X - L(X \mid 1, \mathbf{Z})$ is uncorrelated with $\mathbf{Z}$ by construction. If $\mathbf{Z}$ is a function of $R$ (e.g., size = $\ln R$, or fixed effects that partition the support of $R$), then $L(X \mid 1, \mathbf{Z})$ is a function of $R$. Write $\tilde{X} = \phi(R) + \eta_X - L(\phi(R) + \eta_X \mid 1, \mathbf{Z})$. Under the assumption that $\eta_X$ is mean-zero and uncorrelated with $R$ and with $\mathbf{Z}$ (when $\mathbf{Z}$ is a function of $R$, this holds if $\eta_X \perp R$), we have $L(X \mid 1, \mathbf{Z}) = L(\phi(R) \mid 1, \mathbf{Z}) + 0$, so $\tilde{X} = \phi(R) - L(\phi(R) \mid 1, \mathbf{Z}) + \eta_X$. Thus $\tilde{X}$ is a function of $R$ and $\eta_X$ only. Similarly, the part of $Y$ that is correlated with $\tilde{X}$ is $\psi(R)$ (plus possibly $\eta_Y$ if $\eta_Y$ is correlated with $\eta_X$; we assume they are uncorrelated). Hence $\mathrm{Cov}(Y, \tilde{X}) = \mathrm{Cov}(\psi(R) + \eta_Y, \tilde{X}) = \mathrm{Cov}(\psi(R), \tilde{X}) + \mathrm{Cov}(\eta_Y, \tilde{X})$. If $\eta_Y$ is uncorrelated with $\eta_X$ and with $R$, then $\mathrm{Cov}(\eta_Y, \tilde{X}) = 0$, so $\mathrm{Cov}(Y, \tilde{X}) = \mathrm{Cov}(\psi(R), \tilde{X})$, which is determined only by the distribution of $R$ and the functions $\phi$, $\psi$. Therefore $\beta = \mathrm{Cov}(\psi(R), \tilde{X}) / \mathrm{Var}(\tilde{X})$ depends only on the same-source structure $(R, \phi, \psi)$ and the controls $\mathbf{Z}$, not on the relationship between $D^*$ and $Y^*$.

\textbf{Step 3.} By assumption (ii), $\mathbb{E}[Y^* \mid D^*, R] = \mathbb{E}[Y^* \mid R]$, so the causal effect of $D^*$ on $Y^*$ (given $R$) is zero. The population coefficient $\beta$ in the regression of the proxy $Y$ on the proxy $X$ is a function of $\mathrm{Cov}(Y, \tilde{X})$ and $\mathrm{Var}(\tilde{X})$, which we have shown are determined by $(R, \phi, \psi, \mathbf{Z})$. They do not depend on $D^*$ or $Y^*$. Hence $\beta$ does not identify the causal effect of $D^*$ on $Y^*$; it identifies a reduced-form association induced by the joint dependence of $X$ and $Y$ on $R$.

\textbf{Step 4.} To show that $\beta$ can be nonzero and the $t$-statistic large: suppose $\phi(R)$ and $\psi(R)$ are not constant and are positively or negatively correlated under the distribution of $R$ (e.g., $X$ = rank of $R$, $Y$ = growth of $R$). Then $\mathrm{Cov}(\phi(R), \psi(R)) \neq 0$, and under the orthogonality assumptions above, $\mathrm{Cov}(Y, \tilde{X}) \neq 0$ and $\mathrm{Var}(\tilde{X}) > 0$, so $\beta \neq 0$. In a finite sample of size $n$, the OLS estimator $\hat{\beta}$ converges in probability to $\beta$, and $\sqrt{n}(\hat{\beta} - \beta)$ is asymptotically normal with variance $\sigma^2_\varepsilon / \mathrm{Var}(\tilde{X})$ under standard conditions, where $\sigma^2_\varepsilon$ is the variance of the regression error. Thus the $t$-statistic $t = \hat{\beta} / \widehat{\mathrm{SE}}(\hat{\beta})$ can be large in absolute value (e.g., $|t| > 2$) when $\beta \neq 0$ and $n$ is moderate, even though the true causal effect of $D^*$ on $Y^*$ is zero. \hfill $\square$

\begin{quote}
\textbf{Corollary 1 (Monotone-source proxies: directional and non-vanishing bias).} \textit{Suppose the conditions of Proposition 1 hold. In addition, assume that $X = \phi(R) + \eta_X$ and $Y = \psi(R) + \eta_Y$, where $\phi$ and $\psi$ are non-constant monotone (measurable) functions on the support of $R$, and $\eta_X$, $\eta_Y$ are mean-zero noises with $\mathbb{E}[\eta_X \mid R] = \mathbb{E}[\eta_Y \mid R] = 0$ and $\mathbb{E}[\eta_X \eta_Y \mid R] = 0$. Let $\tilde{X} = X - L(X \mid 1, \mathbf{Z})$. Then: (1) If $\phi$ and $\psi$ are monotone in the same direction, then $\mathrm{Cov}(\psi(R), \tilde{X}) \geq 0$, with strict inequality whenever $\psi(R)$ is not a.s.\ constant after removing the linear projection on $(1, \mathbf{Z})$. (2) If $\phi$ and $\psi$ are monotone in opposite directions, then $\mathrm{Cov}(\psi(R), \tilde{X}) \leq 0$, with strict inequality under the same condition. (3) The population coefficient $\beta = \mathrm{Cov}(Y, \tilde{X}) / \mathrm{Var}(\tilde{X})$ therefore has a systematic sign determined by the monotone source structure, not by the causal effect of $D^*$ on $Y^*$. (4) If the variation in $R$ is not fully absorbed by $\mathbf{Z}$, then $\mathrm{Var}(\tilde{X}) > 0$ and $\beta \neq 0$ generically; the bias is non-vanishing at the population level.}
\end{quote}

\noindent\textit{Proof.} From Proposition 1, $\beta = \mathrm{Cov}(Y, \tilde{X}) / \mathrm{Var}(\tilde{X})$ and under the orthogonality assumptions $\mathrm{Cov}(Y, \tilde{X}) = \mathrm{Cov}(\psi(R), \tilde{X})$. The residual $\tilde{X} = X - L(X \mid 1, \mathbf{Z})$ satisfies $\tilde{X} = \phi(R) - L(\phi(R) \mid 1, \mathbf{Z}) + \eta_X$ when $\eta_X$ is orthogonal to $(1, \mathbf{Z})$. So $\mathrm{Cov}(\psi(R), \tilde{X}) = \mathrm{Cov}(\psi(R), \phi(R)) - \mathrm{Cov}(\psi(R), L(\phi(R) \mid 1, \mathbf{Z}))$. When $\phi$ and $\psi$ are both non-decreasing (or both non-increasing) in $R$, they are positively quadrant dependent, so $\mathrm{Cov}(\phi(R), \psi(R)) \geq 0$. When $\mathbf{Z}$ is a function of $R$ or when the linear projection $L(\phi(R) \mid 1, \mathbf{Z})$ does not reverse the covariation of $\phi(R)$ with $\psi(R)$, the sign of $\mathrm{Cov}(\psi(R), \tilde{X})$ is the same as that of $\mathrm{Cov}(\phi(R), \psi(R))$ (or the second term is smaller in magnitude), so $\mathrm{Cov}(\psi(R), \tilde{X}) \geq 0$ in the same-direction case. If $\phi$ and $\psi$ are monotone in opposite directions, they are negatively quadrant dependent, so $\mathrm{Cov}(\phi(R), \psi(R)) \leq 0$, and under the same condition $\mathrm{Cov}(\psi(R), \tilde{X}) \leq 0$. Strict inequality holds when $\psi(R)$ and $\tilde{X}$ are not a.s.\ constant. Thus (1)--(2) hold. Part (3) follows because $\beta$ has the same sign as $\mathrm{Cov}(\psi(R), \tilde{X})$. For (4), if $\mathbf{Z}$ does not fully absorb the variation in $R$, then $\mathrm{Var}(\tilde{X}) > 0$ and $\mathrm{Cov}(\psi(R), \tilde{X})$ is generically nonzero for non-constant monotone $\phi$, $\psi$. \hfill $\square$

\noindent\textit{Interpretation.} When both the regressor and the outcome are monotone transformations of the same source (e.g., revenue rank and revenue growth), the sign of $\beta$ is mechanically determined by whether the two transformations are in the same or opposite direction. A negative coefficient need not imply a negative causal effect; it may reflect opposite-direction monotone dependence induced by the common source. This rationalizes the negative coefficients in our empirical illustration (Table~\ref{tab:regression}).

\begin{quote}
\textbf{Theorem 2 (Spurious significance under same-source proxy regressions).} \textit{Suppose the conditions of Proposition 1 hold and $\beta \neq 0$ at the population level (e.g., under the conditions of Corollary 1 with non-constant monotone $\phi$, $\psi$ and variation in $R$ not fully absorbed by $\mathbf{Z}$). Consider the panel sample $\{(Y_{it}, X_{i,t-1}, \mathbf{Z}_{it})\}$, $i = 1, \ldots, N$, $t = 1, \ldots, T$, with $N T \to \infty$. Under standard regularity conditions for OLS (finite second moments, appropriate rank condition, and conditions ensuring consistency and asymptotic normality of the OLS estimator with cluster- or heteroskedasticity-robust standard errors), the OLS estimator $\hat{\beta}$ satisfies $\hat{\beta} \xrightarrow{p} \beta$ and $\sqrt{NT}(\hat{\beta} - \beta) \xrightarrow{d} \mathcal{N}(0, V)$ for some $V \in (0, \infty)$. Hence the $t$-statistic $t(\hat{\beta}) = \hat{\beta} / \widehat{\mathrm{SE}}(\hat{\beta})$ satisfies $|t(\hat{\beta})| \xrightarrow{p} +\infty$ as $N T \to \infty$. In particular, even when the true causal effect of $D^*$ on $Y^*$ is zero, the proxy regression can produce asymptotically significant coefficients purely because the proxies share the same underlying source variable.}
\end{quote}

\noindent\textit{Proof.} The population coefficient $\beta$ is the coefficient on $X$ in the linear projection of $Y$ onto $(1, X, \mathbf{Z})$. By assumption $\beta \neq 0$. Under standard OLS regularity conditions (e.g., $\mathbb{E}[\varepsilon_{it} \mid X_{i,t-1}, \mathbf{Z}_{it}] = 0$, finite fourth moments, and non-degeneracy of the design so that the probability limit of $(NT)^{-1} \sum_{i,t} \tilde{X}_{it}^2$ is positive), the OLS estimator $\hat{\beta}$ is consistent: $\hat{\beta} \xrightarrow{p} \beta$. $\sqrt{NT}(\hat{\beta} - \beta) = (NT)^{-1/2} \sum_{i,t} \tilde{X}_{it} \varepsilon_{it} / \widehat{Q} + o_p(1)$, where $\widehat{Q}$ is the sample variance of the residualized regressor, which converges in probability to a positive limit. Under suitable mixing or clustering assumptions, $(NT)^{-1/2} \sum_{i,t} \tilde{X}_{it} \varepsilon_{it}$ is asymptotically normal with mean zero and finite variance, so $\sqrt{NT}(\hat{\beta} - \beta) \xrightarrow{d} \mathcal{N}(0, V)$ for some $V \in (0, \infty)$. The standard error $\widehat{\mathrm{SE}}(\hat{\beta})$ typically satisfies $\widehat{\mathrm{SE}}(\hat{\beta}) = O_p((NT)^{-1/2})$. Thus $t(\hat{\beta}) = \hat{\beta} / \widehat{\mathrm{SE}}(\hat{\beta}) = \beta / \widehat{\mathrm{SE}}(\hat{\beta}) + o_p(1)$. Since $\beta \neq 0$ and $\widehat{\mathrm{SE}}(\hat{\beta}) \xrightarrow{p} 0$ at rate $(NT)^{-1/2}$, we have $|t(\hat{\beta})| \to +\infty$ in probability. \hfill $\square$

\noindent\textit{Interpretation.} Same-source proxy regressions are biased and can become systematically more statistically significant as the sample grows. Thus statistical significance does not rescue proxy validity; larger samples can make structurally induced false positives appear more convincing. Theorem~2 concerns asymptotic behavior ($NT \to \infty$); the controlled experiment in Section~\ref{sec:synthetic-experiment} illustrates the same mechanism in a finite sample ($N=500$, $T=4$).

\begin{quote}
\textbf{Corollary 2 (Residual same-source bias under partial controls).} \textit{Even when controls $\mathbf{Z}$ include transformations of the same source variable $R$ (e.g., $\ln R$, fixed effects, or other covariates built from the same statements), same-source bias remains whenever $X$ and $Y$ retain variation driven by $R$ that is not fully projected out by $\mathbf{Z}$. In particular, controlling for size (e.g., $\ln R$) or for fixed effects that partition the support of $R$ does not eliminate the bias if the residualized regressor $\tilde{X}$ and the outcome $Y$ still covary through the residual part of $R$ in $\tilde{X}$.}
\end{quote}

\noindent\textit{Proof.} This follows directly from Proposition 1: $\beta = \mathrm{Cov}(Y, \tilde{X}) / \mathrm{Var}(\tilde{X})$ where $\tilde{X}$ is the residual from projecting $X$ onto $(1, \mathbf{Z})$. If $\mathbf{Z}$ includes functions of $R$ but does not span the part of $\phi(R)$ that covaries with $\psi(R)$, then $\mathrm{Cov}(\psi(R), \tilde{X}) \neq 0$ and $\beta \neq 0$. Controlling for $\ln R$ or for indicators does not generally absorb all variation in $\phi(R)$ and $\psi(R)$ that is shared, so the bias is residual. \hfill $\square$

\subsubsection{Diagnostic procedure}

Proposition 1 implies that a significant coefficient in a proxy-on-proxy regression is not sufficient evidence of a causal relationship. We propose a \textbf{same-source diagnostic} that applied researchers can implement:

\textbf{Step 1 (Identify same-source component).} For each variable in the regression (regressor, outcome, key controls), note whether it is constructed from the same underlying source (e.g., revenue, accounting statements). If the regressor and at least one outcome are both revenue-derived, same-source bias is possible.

\textbf{Step 2 (Placebo-outcome check).} Construct a \emph{placebo outcome} $\tilde{Y}$ that (a) is derived from the same source (e.g., revenue) so it has similar marginal structure, and (b) is theoretically unrelated to the construct of interest (e.g., a different revenue-based measure that should not be caused by digitalization). Run the same specification with $\tilde{Y}$ as the dependent variable. If the coefficient on $X$ is statistically significant, same-source structure is sufficient to produce significance; the original result cannot be interpreted as evidence of a causal effect of the construct on the outcome.

\textbf{Step 3 (Controlled-experiment check, when feasible).} Under a DGP where the target construct is by design independent of the source variable $R$, generate data and run the same proxy regression. If the coefficient is significant, this demonstrates that same-source structure alone can produce the observed pattern. Our controlled experiment (Section~\ref{sec:synthetic-experiment}) implements this check: true digitalization $\perp$ revenue, yet the proxy regression yields $\hat{\beta} \neq 0$ with large $|t|$.

A study that fails the diagnostic (significant coefficient in the placebo or controlled-experiment check) should not interpret the main regression as identifying the causal effect of the construct; the blueprint in Section~\ref{sec:target-design} (direct measures, strong instruments) indicates what would be required for identification.

\textbf{Practical checklist.} Researchers may summarize same-source risk by reporting: (i) whether the regressor and main outcome are derived from the same raw variable or statement block; (ii) whether a placebo outcome from the same source yields a significant coefficient in the same specification; (iii) whether a controlled experiment with $D^* \perp R$ reproduces a significant coefficient. A ``yes'' on two or more items should be treated as a high-risk proxy design; inference on the target causal effect is not justified.

\subsubsection{Controlled experiment: same-source structure}
\label{sec:synthetic-experiment}

\textbf{Declaration (stated here once):} The data used in this experiment are synthetically generated; this is the only such experiment in the paper, and we do not repeat the term ``synthetic'' in the results or discussion. The experiment is for methodological illustration only.

\textbf{Purpose:} To show that when the data-generating process (DGP) has latent ``true digitalization'' \emph{uncorrelated} with revenue, the same proxy operationalization (industry-year revenue rank, revenue growth) still yields a statistically significant regression coefficient, so the association is driven by same-source (revenue) structure, not by the theoretical relationship. We use industry-year revenue \emph{rank} (percentile) rather than revenue level to mirror the proxy commonly used in the literature (size-based or revenue-based rankings within industry-year), so the experiment is directly comparable to typical applications. \textbf{DGP:} We use a balanced firm-year panel with $N = 500$ firms and $T = 4$ years. For each firm, (i) log revenue follows a random walk, a standard parsimonious assumption for firm revenue dynamics in short panels (so revenue growth is stationary); (ii) latent ``true digitalization'' is drawn independently from a standard normal and is by construction uncorrelated with revenue. We deliberately keep the DGP minimal (no AR(1), trend, or industry heterogeneity) to isolate the same-source mechanism: the point is that even this parsimonious design suffices to produce a significant proxy regression when both regressor and outcome are revenue-based. Richer DGPs (e.g., AR(1) revenue, industry-specific trends) could be explored in future work. We then apply the same variable construction as in the main illustration: industry-year revenue rank (digitalization proxy), lagged one year; revenue growth (environmental performance proxy), clipped to $[-2, 2]$. \textbf{Parameters:} The scale of revenue and growth is chosen so that the means and standard deviations of the proxy variables are in a range consistent with typical firm-level panels (e.g., revenue growth mean near zero, SD on the order of 0.1; revenue rank mean 0.5, SD about 0.29). \textbf{Reproducibility:} We use a fixed random seed (42) and report validation statistics: mean and SD of the environmental performance proxy and lagged digitalization proxy, and the correlation between latent true digitalization and revenue (which should be near zero). Data and code will be deposited in a public repository upon acceptance.
